\conclusion

В результате выполнения выпускной квалификационной работы бакалавра был разработан программный комплекс автоматического распознавания визуальных признаков детских рисунков и сопоставления их с психологическими характеристиками ребёнка нейросетевыми методами. Были решены следующие задачи:

\begin{itemize}
    \item проведён анализ предметной области рисуночных тестов и существующих подходов к автоматизации их интерпретации;
    \item разработан собственный алгоритм предобработки изображений, реализующий локализацию фигуры человека на листе на основе классических операций обработки изображений (адаптивная бинаризация, морфологические операции замыкания, выделение связных компонент, фильтрация и объединение по якорной стратегии);
    \item спроектирована и обучена двухпотоковая собственная свёрточная нейросетевая модель (около 5,5\,млн параметров) для одновременной классификации 117 признаков рисунка;
    \item реализована стратегия обучения, устойчивая к сильному дисбалансу классов: взвешенная перекрёстная энтропия с насыщением весов классов, label smoothing для категориальных признаков, взвешенное случайное сэмплирование, групповое стратифицированное разбиение выборки;
    \item составлен словарь сопоставления 117 технических признаков с 26 каноническими психологическими чертами на основе методики А.~Л.~Венгера;
    \item реализован модуль генерации текстового психологического портрета с ранжированием черт по числу подтверждающих признаков;
    \item разработан микросервисный веб-сервис, объединяющий все модули в единый конвейер обработки с пользовательским интерфейсом, доступным через браузер;
    \item проведено экспериментальное исследование качества модели на тестовой выборке и функциональное тестирование веб-сервиса.
\end{itemize}

На тестовой выборке разработанная система достигла следующих показателей качества: F1\textsubscript{micro} = 0,901, F1\textsubscript{weighted} = 0,875, F1\textsubscript{macro} = 0,521. Значения F1\textsubscript{micro} и F1\textsubscript{weighted} означают, что модель корректно классифицирует около 90\,\% всех предсказаний по 117 признакам в совокупности. Заметно более низкое значение F1\textsubscript{macro} объясняется наличием в выборке нескольких признаков с экстремально малым числом положительных примеров (10 из 117 признаков имеют менее десяти положительных примеров в обучающей выборке), для которых надёжная классификация минорного класса принципиально невозможна без увеличения объёма данных.

Архитектурные особенности разработанного решения --- двухпотоковая модель, объединяющая визуальные признаки изображения и геометрические метаданные bounding box; раздельные проекционные головы для признаков различной сложности обучения; стратегия взвешенных потерь и взвешенного сэмплирования для обработки дисбаланса классов; собственный алгоритм локализации фигуры на основе классических методов обработки изображений --- обеспечивают качество классификации, достаточное для практического применения системы.

Для практического использования разработанные модули объединены в микросервисный веб-сервис, развёртываемый через Docker Compose. Архитектура включает шесть сервисов (предобработка, извлечение признаков, инференс, интерпретация, формирование отчёта, API-шлюз), базу данных PostgreSQL и объектное хранилище MinIO. Веб-интерфейс позволяет загрузить рисунок и получить психологический портрет за 2--3 секунды на CPU без необходимости установки локального окружения.

Программный комплекс автоматизирует рутинную часть обработки рисуночных тестов и позволяет психологу сосредоточиться на тех рисунках, которые требуют детального индивидуального анализа. Это особенно ценно при работе с большими коллекциями рисунков (групповое тестирование в детских садах, школах, исследовательские проекты).

Возможные направления развития работы включают:

\begin{itemize}
    \item расширение датасета редкими классами через целевой сбор рисунков с конкретными признаками, что повысит качество классификации редко встречающихся категорий;
    \item обучение отдельных специализированных моделей для самых сложных признаков по принципу per-criterion fine-tuning поверх общего бэкбона;
    \item применение более крупных архитектур и техник self-supervised предобучения на расширенной коллекции неразмеченных рисунков;
    \item расширение комплекса на другие рисуночные тесты (тесты <<Дом-Дерево-Человек>>, <<Несуществующее животное>>, <<Моя семья>>).
\end{itemize}

Исходный код разработанного программного комплекса размещён в открытом репозитории. QR-код для быстрого перехода к репозиторию приведён в приложении А на рисунке~\ref{fig:repo-qr}.

Полученный программный комплекс представляет собой цельное решение, в основе которого лежат современные методы глубокого обучения, классические алгоритмы обработки изображений, корректная инженерная работа с реальным несбалансированным датасетом и интеграция доменных психологических знаний. Результаты работы могут быть использованы как основа для дальнейших исследований и практических применений автоматизированной психологической диагностики на основе рисуночных тестов.
